\documentclass[11pt]{article}
\usepackage{fontspec}
\setmainfont{Latin Modern Roman}
\usepackage{amsmath,amssymb,mathtools,unicode-math}
\setmathfont{Latin Modern Math}
\usepackage{xurl,hyperref}
\hypersetup{hidelinks,pdftitle={Problem statement}}
\newcommand{\sref}[2]{\hyperlink{source:#1:#2}{[S#2]}}
\newcommand{\cataloguescope}[1]{\paragraph{Mathematical question.}#1}
\begin{document}
\section{Hamilton cycle in small d-diregular graphs}\label{statement}

\subsection{Definitions and mathematical statement}

\paragraph{Shared notation.}$\mathbb N=\{1,2,\ldots\}$, and $\mathbb Z,\mathbb Q,\mathbb R,\mathbb C$ denote the integers, rationals, reals and complex numbers. Any other conventions are specified locally.


An oriented graph has no loops and at most one directed arc between each unordered pair of vertices. Write $d^+(v),d^-(v)$ for outdegree and indegree. The source calls it $d$-diregular when both are at least $d$ at every vertex. A $d$-regular oriented graph has both equal to $d$. Strong $2$-connectivity means that deleting any one vertex leaves a strongly connected digraph. A Hamilton cycle is a directed cycle through every vertex.
\paragraph{Statement recorded in the supplied source.}
For every integer $d>2$, does every $d$-diregular oriented graph of order at most $4d+1$ have a Hamilton cycle? The source also asks: for every integer $d>2$, does every $d$-regular strongly $2$-connected oriented graph on at most $6d$ vertices have a Hamilton cycle? \sref{statement}{2}
\subsection{Short English statement}
Do the stated order and degree conditions force directed Hamilton cycles, with a larger order bound under strong 2-connectivity and exact regularity?
\subsection{Sources}
\begin{description}
\item[\hypertarget{source:statement:1}{S1}] ULAM.AI and the UnsolvedMath Contributors, UnsolvedMath: A Curated Collection of Open Mathematics Problems, catalogue entry. CC BY 4.0. \url{https://huggingface.co/datasets/ulamai/UnsolvedMath}.
\item[\hypertarget{source:statement:2}{S2}] Open Problem Garden, Hamilton cycle in small d-diregular graphs, original node 47028, statement and mathematical discussion. \url{https://www.openproblemgarden.org/op/hamilton_cycle_in_small_d_diregular_graphs}.
\end{description}
\label{end:statement}
\end{document}
